\subsection{The cominuscule representation of $G$ associated with $M$}
\subsubsection{The representation}
\begin{definition}
A simple root $\beta\in\Pi$, is \emph{cominuscule} w.r.t. the root system $R$ if $n_\beta(\alpha)\in\{-1,0,1\}$ for all $\alpha\in R$, or equivalently if $n_\beta(\Theta)=1$, where $\Theta$ is the highest root of $R$.

An irreducible representation of $G$ whose highest weight is equal to $\varpi_\beta$ for $\beta$ a cominuscule root of $R$ is called a \emph{cominuscule representation}.
\end{definition}
The following is well-known, but we include a proof.
\begin{proposition}
There is a unique simple noncompact root. This root is long and cominuscule.
\end{proposition}
\begin{notation}
This unique simple noncompact root will be denoted by $\zeta$.
\end{notation}
\begin{proof}
Since $R$ is irreducible, the highest root $\Theta$ of $R$ is unique and $\Theta$ has a positive coefficient on every simple root. Now, if there are more than one simple noncompact root, or if $n_\zeta(\Theta)>1$ for some noncompact simple root $\zeta$, then $\langle\Theta,z\rangle>2$. Contradiction. Hence $\zeta$ is unique and cominuscule. Assume that $\zeta$ is short (so that $R$ is not simply laced). Then, since $\Theta$ is long,
\[
\langle\Theta,\zeta^\jepPubCoroot\rangle=\langle\zeta,\Theta^\jepPubCoroot\rangle\frac{(\text{root length of }\Theta)^2}{(\text{root length of }\zeta)^2}\geqslant2.
\]
Since $s_\zeta(\Theta)$ is a positive root, this implies that $n_\zeta(\Theta)\geqslant2$, contradicting the fact that $\zeta$ is cominuscule.
\end{proof}

Note that the maximal parabolic subalgebra $\mathfrak p=\mathfrak k\oplus\mathfrak m^+$ of $\mathfrak g$ is the standard parabolic defined by the root $\zeta$:
\[
\mathfrak p=\mathfrak t\oplus\jepPubOplus_{\alpha:n_\zeta(\alpha)\geqslant0}\mathfrak g_\alpha.
\]
Recall that $P$ is the corresponding parabolic subgroup of $G$ and that the compact dual of $M$ is $\check M=G/P$. By \cite{jep93Mur59}, since $G$ is simply connected, the Picard group $\operatorname{Pic}(G/P)$ is isomorphic to the group of characters $\jepPubScript{X}(P)$ of $P$. Since $\mathfrak p=\mathfrak z\oplus[\mathfrak p,\mathfrak p]$, $\jepPubScript{X}(P)$ is isomorphic to $\mathbb Z$ and thus it is generated by the smallest positive character of $P$, namely $\varpi:=\varpi_\zeta$. Here we embed $\jepPubScript{X}(P)$ in $\mathfrak t^*$ via $\lambda\mapsto d\lambda$, the differential of $\lambda$ at the unit element of $T$. Moreover, the isomorphism $\iota:\jepPubScript{X}(P)\simeq\operatorname{Pic}(G/P)$ is given by $\lambda\mapsto(G\times\mathbb C_\lambda)/P$, where $\mathbb C_\lambda$ denotes the $1$-dimensional $P$-module defined by $\lambda$. In particular, $\jepPubScript{L}:=\iota(\varpi)$ is a generator of $\operatorname{Pic}(\check M)$.

Let now $\jepPubBbb{E}$ be the cominuscule irreducible representation of $G$ whose highest weight is $\varpi=\varpi_\zeta$. Let $\jepPubBbb{E}_\varpi$ be the $\varpi$-eigenspace of $\jepPubBbb{E}$. Then $\jepPubBbb{E}_\varpi$ is 1-dimensional and its stabilizer in $G$ is $P$. This gives a $G$-equivariant holomorphic and isometric embedding of $\check M=G/P$ in the projective space $\mathbb P\jepPubBbb{E}$. It is called the \emph{first canonical embedding} of $\check M$, and since $\jepPubBbb{E}_\varpi\simeq\mathbb C_\varpi$ as $P$-modules, we have $\jepPubScript{L}\simeq\jepPubScript{O}_{\mathbb P\jepPubBbb{E}}(-1)$. See e.g. \cite{jep93NT76} for more details.

\begin{remark}
The complex group $G$ may have several cominuscule simple roots or representations. In fact, cominuscule simple roots are in one to one correspondence with Hermitian real forms $G_{\mathbb R}$ (or $\check G_{\mathbb R}$) of $G$. When we will need a case by case argument, we will rather use the classification of the cominuscule roots than the classification of the real groups $G_{\mathbb R}$. This correspondence is shown in the first two columns of Table~1.
\end{remark}
We now begin our study of the cominuscule representation $\jepPubBbb{E}$ of $G$.
\begin{notation}
We denote by $\mu_0$ the lowest weight of $\jepPubBbb{E}$ and by $X(\jepPubBbb{E})$ the set of weights of $\jepPubBbb{E}$. For $\chi\in X(\jepPubBbb{E})$, we write $\jepPubBbb{E}_\chi$ for the corresponding weight space. Recall that $\varpi$ is the highest weight of $\jepPubBbb{E}$.
\end{notation}
The fact that $\jepPubBbb{E}$ is cominuscule has the following consequence on the weights of $\jepPubBbb{E}$.
\begin{lemma}
For any weight $\chi$ of $\jepPubBbb{E}$, and any root $\alpha\in R$, $\langle\chi,\alpha^\jepPubCoroot\rangle\in\{-2,-1,0,1,2\}$, and the equality $\langle\chi,\alpha^\jepPubCoroot\rangle=\pm2$ implies that $\alpha$ is short (hence that $R$ is not simply laced).
\end{lemma}
\begin{proof}
For the highest weight $\varpi$ of $\jepPubBbb{E}$, the results follows from the facts that $\langle\varpi,\alpha^\jepPubCoroot\rangle=n_\zeta(\alpha)\langle\zeta,\alpha^\jepPubCoroot\rangle/\langle\alpha,\zeta^\jepPubCoroot\rangle$ and $\zeta$ is a long root. The result still holds if $\varpi$ is replaced by $w\cdot\varpi$, where $w\in W$ is arbitrary, and since any weight of $\jepPubBbb{E}$ is in the convex hull of $W\cdot\varpi$, it holds for any weight.
\end{proof}
We deduce that the structure of $\jepPubBbb{E}$ with respect to the action of $\mathfrak g_\alpha$ for $\alpha$ a long root is particularly simple.
\begin{lemma}
Let $\alpha$ be a long root and let $\chi$ be a weight of $\jepPubBbb{E}$. We have
\[
\mathfrak g_{-\alpha}\cdot \jepPubBbb{E}_\chi=\begin{cases}\jepPubBbb{E}_{\chi-\alpha}&\text{if }\langle\chi,\alpha^\jepPubCoroot\rangle=1,\\\{0\}&\text{otherwise}\end{cases}.
\]
\end{lemma}
\begin{proof}
Let $\alpha$ be long and let $\mathfrak{sl}(\alpha)$ be the Lie subalgebra of $\mathfrak g$ isomorphic to $\mathfrak{sl}_2$ corresponding to $\alpha$. Let $S\subset \jepPubBbb{E}$ be the $\mathfrak{sl}(\alpha)$-submodule generated by $\jepPubBbb{E}_\chi$. By Lemma~2.8 and the representation theory of $\mathfrak{sl}_2$, any irreducible component $V$ of $S$ is an $\mathfrak{sl}(\alpha)$-module of dimension $1$ or $2$.

We therefore have only three possibilities. The first case is when $V=V_\chi$, $\mathfrak g_\alpha\cdot V_\chi=\{0\}$ and $\mathfrak g_{-\alpha}\cdot V_\chi=\{0\}$. In this case, $\langle\chi,\alpha^\jepPubCoroot\rangle=0$. The second case is when $V=V_\chi\oplus V_{\chi-\alpha}$, $\mathfrak g_\alpha\cdot V_\chi=\{0\}$ and $\mathfrak g_{-\alpha}\cdot V_\chi=V_{\chi-\alpha}$. In this case, $\langle\chi,\alpha^\jepPubCoroot\rangle=1$ (and $\langle\chi-\alpha,\alpha^\jepPubCoroot\rangle=-1$). The third (symmetric) case is when $V=V_\chi\oplus V_{\chi+\alpha}$, $\mathfrak g_\alpha\cdot V_\chi=V_{\chi+\alpha}$ and $\mathfrak g_{-\alpha}\cdot V_\chi=\{0\}$. In this case, $\langle\chi,\alpha^\jepPubCoroot\rangle=-1$ (and $\langle\chi+\alpha,\alpha^\jepPubCoroot\rangle=1$).

If $\langle\chi,\alpha^\jepPubCoroot\rangle=1$, we deduce that $S=S_\chi\oplus S_{\chi-\alpha}$ and that $\mathfrak g_{-\alpha}\cdot S_\chi=S_{\chi-\alpha}$. We have $s_\alpha(\chi)=\chi-\alpha$ so $\dim(\jepPubBbb{E}_\chi)=\dim(\jepPubBbb{E}_{\chi-\alpha})$. The lemma is proved in this case. If $\langle\chi,\alpha^\jepPubCoroot\rangle\leqslant0$, we see that $\mathfrak g_{-\alpha}\cdot \jepPubBbb{E}_\chi=\{0\}$.
\end{proof}

\subsubsection{The grading}
\begin{notation}
Let $z_{\mathtt{max}}$ denote the number $\langle\varpi,z\rangle$.
\end{notation}
\begin{remark}
The precise value of $z_{\mathtt{max}}$ will not play any role in the following, however it has already been computed in \cite[p.~214-216]{jep93KM10}): with our choice of $z\in\mathfrak z$ we have $z_{\mathtt{max}}=2\dim\check M/c_1(\check M)$, where $c_1(\check M)$ is the first Chern number of $\check M$. For example if $G_{\mathbb R}=\mathrm{SU}(p,q)$, $z_{\mathtt{max}}=2pq/(p+q)$.
\end{remark}
\begin{proposition}
The set $\{\langle\chi,z\rangle\mid\chi\in X(\jepPubBbb{E})\}$ is the set
\[
\{z_{\mathtt{max}},z_{\mathtt{max}}-2,\ldots,z_{\mathtt{max}}-2r_M\}.
\]
\end{proposition}
\begin{proof}
It follows from \cite[Th.~2.1]{jep93RRS92} that the $W$-orbit of the weight $\varpi$ is exactly the set of weights of the form $\varpi-\sum_{i=1}^k\delta_i$, where $(\delta_i)_{1\leqslant i\leqslant k}$ is a family of noncompact long strongly orthogonal roots. In fact, the orbit $W\cdot\varpi$ is in bijection with $W/W_P$, so the equivalence between items~(c) and~(e) in the cited theorem proves the claim. For any $i$, we have $\langle\delta_i,z\rangle=2$, thus we have the equality of sets
\[
\{\langle\mu,z\rangle\mid\mu\in W\cdot\varpi\}=\{z_{\mathtt{max}},z_{\mathtt{max}}-2,\ldots,z_{\mathtt{max}}-2r_M\}.
\]
In particular, $\langle\mu_0,z\rangle=z_{\mathtt{max}}-2r_M$ and for $\chi\in X(\jepPubBbb{E})$, we have $z_{\mathtt{max}}-2r_M\leqslant\langle\chi,z\rangle\leqslant z_{\mathtt{max}}$. The result of the proposition now follows from the fact that $2$ is a divisor of $\langle\alpha,z\rangle$ for any root $\alpha$.
\end{proof}
Now we can introduce the grading of $\jepPubBbb{E}$.
\begin{definition}
For a relative integer $i$, let
\[
\jepPubBbb{E}_i:=\jepPubOplus_{\chi:\langle\chi,z\rangle=z_{\mathtt{max}}-2i}\jepPubBbb{E}_\chi.
\]
\end{definition}
This grading corresponds to the decomposition of $\jepPubBbb{E}$ into irreducible $K$-modules.
\begin{proposition}
The $K$-modules $\jepPubBbb{E}_i$ are irreducible.
\end{proposition}
\begin{proof}
This might be well-known to experts, but we include a proof for completeness. We give a case by case argument and use Table~1.

In type $A_{n-1}$, we have $\jepPubBbb{E}=\wedge^{r_M}(\mathbb C^{r_M}\oplus\mathbb C^{n-r_M})$ and thus $\jepPubBbb{E}_i=\wedge^{r_M-i}\mathbb C^{r_M}\otimes\wedge^i\mathbb C^{n-r_M}$: this is an irreducible $\mathrm S(\mathrm{GL}_{r_M}\times\mathrm{GL}_{n-r_M})$-module.

In type $B_n$, we have $\jepPubBbb{E}=\mathbb C^{2n+1}=\mathbb C\oplus\mathbb C^{2n-1}\oplus\mathbb C$, and each summand is an irreducible $\mathrm{Spin}_{2n-1}$-module, hence an irreducible $K$-module. In type $(D_n,\varpi_1)$ the situation is similar.

In type $(D_n,\varpi_n)$, $\jepPubBbb{E}$ is the spinor representation of the spin group and, according to \cite{jep93Che97}, we have $\jepPubBbb{E}=\wedge^0\mathbb C^n\oplus\wedge^2\mathbb C^n\oplus\cdots\oplus\wedge^{2p}\mathbb C^n$ (where $p=\lfloor n/2\rfloor$). Thus $\jepPubBbb{E}_i=\wedge^{2i}\mathbb C^n$, and this is an irreducible $\mathrm{GL}_n$-module.

For the types $E_6$ and $E_7$, we use models of these exceptional Lie algebras and their minuscule representations, as given for example in \cite{jep93Man06}. In type $E_6$, we have $\jepPubBbb{E}=\mathbb C\oplus V_{16}\oplus V_{10}$, where $V_{16}$ is a spinor representation and $V_{10}$ the vector representation of the spin group $\mathrm{Spin}_{10}$. In type $E_7$, we have $\jepPubBbb{E}=\mathbb C\oplus V_{27}\oplus V'_{27}\oplus\mathbb C$, where $V_{27}$ and $V'_{27}$ are the two minuscule representations of a group of type $E_6$. In both cases, the $K$-modules $\jepPubBbb{E}_i$ are irreducible.

We now deal with the case of type $C_n$. The following representation-theoretic argument has been suggested to us by an anonymous referee that we would like to thank. We denote by $\mathbb C^{2n}=\mathbb C_a^n\oplus\mathbb C_b^n$ a symplectic $2n$-dimensional space, with $\mathbb C_a^n$ and $\mathbb C_b^n$ supplementary isotropic subspaces. We then have $\jepPubBbb{E}=\bigl(\wedge^n\mathbb C^{2n}\bigr)_\omega$, where the symbol $\omega$ means that we take in $\wedge^n\mathbb C^{2n}$ the irreducible $\mathrm{Sp}_{2n}$-submodule containing the highest weight line $\wedge^n\mathbb C_a^n$. More precisely, by \cite[Th.~17.5]{jep93FH91}, we have a decomposition of $\wedge^n\mathbb C^{2n}$ as an $\mathrm{Sp}_{2n}$-module as follows: $\wedge^n\mathbb C^{2n}=\jepPubBbb{E}\oplus\omega\wedge(\wedge^{n-2}\mathbb C^{2n})$.

Thus, $\wedge^{n-i}\mathbb C_a^n\otimes\wedge^i\mathbb C_b^n=\jepPubBbb{E}_i\oplus\omega\wedge(\wedge^{n-i-1}\mathbb C_a^n)\otimes\wedge^{i-1}\mathbb C_b^n.$ We now consider this as a representation of $K=\mathrm{GL}_n$. As $K$-modules, we have $\mathbb C_b^n\simeq(\mathbb C_a^n)^\star$. We thus get $\wedge^{n-i}\mathbb C_a^n\otimes\wedge^{n-i}\mathbb C_a^n\simeq\jepPubBbb{E}_i\oplus\wedge^{n-i-1}\mathbb C_a^n\otimes\wedge^{n-i+1}\mathbb C_a^n.$ From this last equation, it follows that $\jepPubBbb{E}_i$ is the Cartan square of $\wedge^{n-i}\mathbb C_a^n$ and therefore it is an irreducible $\mathrm{GL}_n$-module.
\end{proof}
\begin{proposition}
We have the following properties:
\begin{enumerate}[label=(\alph*)]
\item $\jepPubBbb{E}=\jepPubBbb{E}_0\oplus\jepPubBbb{E}_1\oplus\cdots\oplus\jepPubBbb{E}_{r_M}$
\item $\jepPubBbb{E}_0=\jepPubBbb{E}_\varpi$
\item $\jepPubBbb{E}_{i+1}=\mathfrak m^-\cdot\jepPubBbb{E}_i$
\item The map $\jepPubBbb{E}_0\otimes\mathfrak m^-\to\jepPubBbb{E}_1$ is an isomorphism
\end{enumerate}
\end{proposition}
\begin{proof}
Only the last two points need a proof. Let $U(\mathfrak m^-)$ denote the enveloping algebra of $\mathfrak m^-$. The third point follows from the fact that $\jepPubBbb{E}=U(\mathfrak m^-)\cdot\jepPubBbb{E}_\varpi$ and the fact that for $\alpha$ a root of $\mathfrak m^-$, we have $\langle\alpha,z\rangle=-2$. The last point follows by Schur’s lemma since $\mathfrak m^-$ and $\jepPubBbb{E}_1$ are irreducible $\mathfrak k$-modules and $\jepPubBbb{E}_0$ is 1-dimensional.
\end{proof}
\begin{remark}
Those statements are well-known. In the case where $\jepPubBbb{E}$ is of tube type, they are proved in \cite[Prop.~5.2]{jep93Gro94}.
\end{remark}

